% =====================================================
% 题型：四面体体积比选择题 (q_solid_tetrahedron_vol.tex)
% =====================================================
% =====================================================
% 难度：★★ (中档)
% =====================================================
\newcommand{\qSolidTetrahedronVolStem}{如图，在四面体 \(A-BCD\) 中，\(E, F, G\) 分别为线段 \(AB, AC, AD\) 上的点，若平面 \(EFG \parallel\) 平面 \(BCD\)，\(\dfrac{1}{2}\vec{BC} = \vec{AF} - \vec{AE}\)，则四面体 \(A-BCD\) 与三棱锥 \(A-EFG\) 的体积之比为 \hfill （\quad）}

\newcommand{\qSolidTetrahedronVolOptions}{%
  \mcoptions[4]{\(2\)}{\(3\)}{\(4\)}{\(8\)}%
}

\newcommand{\qSolidTetrahedronVolFigure}[1][1.3]{%
  \begin{tikzpicture}[scale=#1, >=stealth]
    \coordinate (A) at (1.0, 2.0);
    \coordinate (B) at (0, 0);
    \coordinate (C) at (1.2, -0.5);
    \coordinate (D) at (2.2, 0.3);
    
    \coordinate (E) at ($(A)!0.5!(B)$);
    \coordinate (F) at ($(A)!0.5!(C)$);
    \coordinate (G) at ($(A)!0.5!(D)$);
    
    % Back face lines (dashed)
    \draw[dashed, thick, gray!80] (B) -- (D);
    \draw[dashed, thick] (E) -- (G);
    
    \ifteacher
      % 教师端添加：比例引导线
      \draw[dashed, semithick, gray] (B) -- (F);
      \draw[dashed, semithick, gray] (C) -- (G);
    \fi
    
    % Front face lines (solid)
    \draw[thick] (A) -- (B) -- (C) -- (D) -- (A) -- (C);
    \draw[very thick] (E) -- (F) -- (G);
    
    \node[above] at (A) {\small $A$};
    \node[left] at (B) {\small $B$};
    \node[below] at (C) {\small $C$};
    \node[right] at (D) {\small $D$};
    \node[left] at (E) {\small $E$};
    \node[below left] at (F) {\small $F$};
    \node[right] at (G) {\small $G$};
  \end{tikzpicture}%
}

\newcommand{\qSolidTetrahedronVolAnswer}{D}

\newcommand{\qSolidTetrahedronVolSolution}{%
  因为 \(\vec{AF} - \vec{AE} = \vec{EF}\)，且已知 \(\dfrac{1}{2}\vec{BC} = \vec{AF} - \vec{AE}\)，所以 \(\vec{EF} = \dfrac{1}{2}\vec{BC}\)。
  因为平面 \(EFG \parallel\) 平面 \(BCD\)，由面面平行的性质可得 \(EF \parallel BC\)，且截面三角形 \(\triangle EFG\) 相似于底面 \(\triangle BCD\)。
  其相似比为 \(\dfrac{EF}{BC} = \dfrac{1}{2}\)。
  由相似多面体的体积之比等于对应线段之比的立方，可得三棱锥 \(A-EFG\) 与四面体 \(A-BCD\) 的体积比为：
  \(\dfrac{V_{A-EFG}}{V_{A-BCD}} = \left(\dfrac{1}{2}\right)^3 = \dfrac{1}{8}\)。
  因此，四面体 \(A-BCD\) 与三棱锥 \(A-EFG\) 的体积之比为 \(8\)。
  故选 D。%
}
